\section{Può tornare utile}

\begin{softnote}
	Questa sezione contiene materiale di riferimento (appendici) e domande d'esame.
\end{softnote}

\subsection{Appendice A (Sommatorie)}

\subsubsection{Formule fondamentali}

\begin{align*}
	\sum_{i=1}^{n} 1 &= n \\
	\sum_{i=1}^{n} i &= \frac{n(n+1)}{2} = \Theta(n^2) \\
	\sum_{i=1}^{n} i^2 &= \frac{n(n+1)(2n+1)}{6} = \Theta(n^3) \\
	\sum_{i=1}^{n} i^k &= \Theta(n^{k+1}) \quad \text{per } k \geq 0 \\
	\sum_{i=0}^{n} x^i &= \frac{x^{n+1}-1}{x-1} \quad \text{per } x \neq 1 \\
	\sum_{i=0}^{\infty} x^i &= \frac{1}{1-x} \quad \text{per } |x| < 1
\end{align*}

\subsubsection{Serie geometrica}

\begin{nota}
	La serie geometrica è fondamentale nell'analisi degli algoritmi ricorsivi.
	
	\textbf{Formula generale:}
	\[
	\sum_{i=0}^{n} ar^i = a \cdot \frac{r^{n+1}-1}{r-1}
	\]
	
	\textbf{Caso particolare importante} ($r = 2$):
	\[
	\sum_{i=0}^{n} 2^i = 2^{n+1} - 1
	\]
	
	\textbf{Serie infinita convergente} ($|r| < 1$):
	\[
	\sum_{i=0}^{\infty} ar^i = \frac{a}{1-r}
	\]
\end{nota}

\subsubsection{Serie armonica}

\[
H_n = \sum_{i=1}^{n} \frac{1}{i} = 1 + \frac{1}{2} + \frac{1}{3} + \cdots + \frac{1}{n} = \Theta(\log n)
\]

\begin{nota}
	La serie armonica cresce \textbf{logaritmicamente}. Utile per analizzare:
	\begin{itemize}
		\item Quicksort (numero medio di confronti)
		\item Hashing (lunghezza catene)
		\item Alberi bilanciati (altezza media)
	\end{itemize}
	
	\textbf{Approssimazione:} $H_n \approx \ln n + \gamma$ dove $\gamma \approx 0.5772$ (costante di Eulero-Mascheroni)
\end{nota}

\subsubsection{Tecniche di manipolazione}

\textbf{1. Separazione dei termini:}
\[
\sum_{i=1}^{n} (a_i + b_i) = \sum_{i=1}^{n} a_i + \sum_{i=1}^{n} b_i
\]

\textbf{2. Estrazione di costanti:}
\[
\sum_{i=1}^{n} c \cdot a_i = c \sum_{i=1}^{n} a_i
\]

\textbf{3. Telescoping (sommatoria telescopica):}
\[
\sum_{i=1}^{n} (a_i - a_{i-1}) = a_n - a_0
\]

\begin{esempio}[Telescoping applicato]
	\[
	\sum_{i=1}^{n} \frac{1}{i(i+1)} = \sum_{i=1}^{n} \left(\frac{1}{i} - \frac{1}{i+1}\right) = 1 - \frac{1}{n+1} = \frac{n}{n+1}
	\]
\end{esempio}

\textbf{4. Cambio di variabile:}

Se $j = n - i$, allora:
\[
\sum_{i=0}^{n} f(i) = \sum_{j=n}^{0} f(n-j) = \sum_{j=0}^{n} f(n-j)
\]

\subsubsection{Approssimazione con integrali}

Per $f$ crescente:
\[
\int_{1}^{n} f(x)\,dx \leq \sum_{i=1}^{n} f(i) \leq f(n) + \int_{1}^{n} f(x)\,dx
\]

\begin{esempio}[Somma di potenze]
	\begin{align*}
		\sum_{i=1}^{n} i^k &\approx \int_{1}^{n} x^k\,dx = \frac{n^{k+1}}{k+1} \\
		&\Rightarrow \sum_{i=1}^{n} i^k = \Theta(n^{k+1})
	\end{align*}
\end{esempio}

\subsection{Appendice B (Insiemi, relazioni, funzioni, grafi, alberi)}

\subsubsection{Insiemi}

\begin{definizione}[Insieme]
	Collezione non ordinata di elementi distinti. Notazione: $A = \{1, 2, 3\}$
\end{definizione}

\textbf{Operazioni fondamentali:}
\begin{itemize}
	\item \textbf{Unione:} $A \cup B = \{x : x \in A \text{ oppure } x \in B\}$
	\item \textbf{Intersezione:} $A \cap B = \{x : x \in A \text{ e } x \in B\}$
	\item \textbf{Differenza:} $A \setminus B = \{x : x \in A \text{ e } x \notin B\}$
	\item \textbf{Complemento:} $\overline{A} = \{x : x \notin A\}$ (rispetto a un universo)
\end{itemize}

\begin{nota}
	\textbf{Cardinalità:} $|A|$ = numero di elementi in $A$
	
	\textbf{Insieme delle parti:} $\mathcal{P}(A) = $ insieme di tutti i sottoinsiemi di $A$
	
	Se $|A| = n$ allora $|\mathcal{P}(A)| = 2^n$
\end{nota}

\subsubsection{Relazioni}

\begin{definizione}[Relazione]
	Una relazione $R$ da $A$ a $B$ è un sottoinsieme di $A \times B$ (prodotto cartesiano).
	
	Scriviamo $aRb$ se $(a,b) \in R$.
\end{definizione}

\textbf{Proprietà di relazioni su $A$ (relazione binaria):}
\begin{itemize}
	\item \textbf{Riflessiva:} $\forall a \in A: aRa$
	\item \textbf{Simmetrica:} $\forall a,b \in A: aRb \Rightarrow bRa$
	\item \textbf{Transitiva:} $\forall a,b,c \in A: (aRb \land bRc) \Rightarrow aRc$
	\item \textbf{Antisimmetrica:} $\forall a,b \in A: (aRb \land bRa) \Rightarrow a = b$
\end{itemize}

\begin{definizione}[Relazione di equivalenza]
	Una relazione è di \textbf{equivalenza} se è riflessiva, simmetrica e transitiva.
	
	Esempio: uguaglianza ($=$), congruenza modulo $n$
\end{definizione}

\begin{definizione}[Ordine parziale]
	Una relazione è di \textbf{ordine parziale} se è riflessiva, antisimmetrica e transitiva.
	
	Esempio: $\leq$ sui numeri reali, $\subseteq$ sugli insiemi
\end{definizione}

\subsubsection{Funzioni}

\begin{definizione}[Funzione]
	Una funzione $f: A \to B$ assegna a ogni elemento di $A$ esattamente un elemento di $B$.
	
	\begin{itemize}
		\item $A$ = \textbf{dominio}
		\item $B$ = \textbf{codominio}
		\item $f(A) = \{f(a) : a \in A\}$ = \textbf{immagine}
	\end{itemize}
\end{definizione}

\textbf{Tipi di funzioni:}
\begin{itemize}
	\item \textbf{Iniettiva} (uno-a-uno): $f(a_1) = f(a_2) \Rightarrow a_1 = a_2$
	\item \textbf{Suriettiva} (sopra): $\forall b \in B, \exists a \in A: f(a) = b$
	\item \textbf{Biettiva} (biunivoca): iniettiva e suriettiva
\end{itemize}

\begin{nota}
	Una funzione biettiva ha \textbf{funzione inversa} $f^{-1}: B \to A$ tale che:
	\[
	f^{-1}(f(a)) = a \quad \text{e} \quad f(f^{-1}(b)) = b
	\]
\end{nota}

\subsubsection{Grafi}

\begin{definizione}[Grafo]
	Un grafo $G = (V, E)$ consiste di:
	\begin{itemize}
		\item $V$ = insieme di \textbf{vertici} (o nodi)
		\item $E \subseteq V \times V$ = insieme di \textbf{archi}
	\end{itemize}
\end{definizione}

\textbf{Tipi di grafi:}
\begin{itemize}
	\item \textbf{Grafo non orientato:} $(u,v) \in E \Leftrightarrow (v,u) \in E$
	\item \textbf{Grafo orientato} (digrafo): gli archi hanno una direzione
	\item \textbf{Grafo pesato:} ogni arco ha un peso $w: E \to \mathbb{R}$
\end{itemize}

\textbf{Terminologia:}
\begin{itemize}
	\item \textbf{Grado} di $v$: numero di archi incidenti
	\begin{itemize}
		\item Grafo orientato: \textit{grado entrante} + \textit{grado uscente}
	\end{itemize}
	\item \textbf{Cammino:} sequenza di vertici $v_1, v_2, \ldots, v_k$ con $(v_i, v_{i+1}) \in E$
	\item \textbf{Ciclo:} cammino con $v_1 = v_k$
	\item \textbf{Connesso:} esiste un cammino tra ogni coppia di vertici
\end{itemize}

\textbf{Rappresentazioni:}
\begin{enumerate}
	\item \textbf{Matrice di adiacenza} ($n \times n$):
	\[
	A[i][j] = \begin{cases}
		1 & \text{se } (i,j) \in E \\
		0 & \text{altrimenti}
	\end{cases}
	\]
	Spazio: $\Theta(n^2)$, accesso arco: $O(1)$
	
	\item \textbf{Liste di adiacenza:}
	Array di $n$ liste, $\text{Adj}[u]$ contiene i vicini di $u$
	
	Spazio: $\Theta(n + m)$, iterare vicini: $O(\deg(u))$
\end{enumerate}

\subsubsection{Alberi}

\begin{definizione}[Albero]
	Un albero è un grafo \textbf{connesso} e \textbf{aciclico}.
	
	Equivalentemente:
	\begin{itemize}
		\item Grafo connesso con $n$ nodi e $n-1$ archi
		\item Esiste un unico cammino tra ogni coppia di nodi
	\end{itemize}
\end{definizione}

\textbf{Albero radicato:} albero con un nodo speciale detto \textbf{radice}.
\begin{itemize}
	\item \textbf{Padre/Figlio:} $u$ è padre di $v$ se $(u,v) \in E$ e $u$ è più vicino alla radice
	\item \textbf{Foglia:} nodo senza figli
	\item \textbf{Altezza:} lunghezza del cammino più lungo dalla radice a una foglia
	\item \textbf{Profondità} di $v$: distanza dalla radice a $v$
\end{itemize}

\begin{definizione}[Albero binario]
	Albero radicato in cui ogni nodo ha al più 2 figli (sinistro e destro).
	
	\textbf{Albero binario completo:} tutte le foglie allo stesso livello, ogni nodo interno ha 2 figli.
	
	Un albero binario completo di altezza $h$ ha $2^{h+1} - 1$ nodi.
\end{definizione}

\textbf{Visite di alberi binari:}
\begin{itemize}
	\item \textbf{Preorder:} radice $\to$ sinistro $\to$ destro
	\item \textbf{Inorder:} sinistro $\to$ radice $\to$ destro
	\item \textbf{Postorder:} sinistro $\to$ destro $\to$ radice
\end{itemize}

\subsection{Appendice C (Calcolo combinatorio e probabilità)}

\subsubsection{Principi di conteggio}

\textbf{1. Principio del prodotto:}

Se un evento può avvenire in $m$ modi e un altro in $n$ modi, allora entrambi possono avvenire in $m \times n$ modi.

\begin{esempio}
	Quante stringhe binarie di lunghezza 3?
	
	Ogni posizione: 2 scelte $\Rightarrow$ totale: $2 \times 2 \times 2 = 8$
\end{esempio}

\textbf{2. Principio della somma:}

Se un evento può avvenire in $m$ modi disgiunti o in $n$ modi disgiunti, può avvenire in $m + n$ modi.

\subsubsection{Permutazioni e combinazioni}

\begin{definizione}[Permutazione]
	Numero di modi per ordinare $n$ oggetti distinti:
	\[
	P(n) = n! = n \times (n-1) \times \cdots \times 2 \times 1
	\]
\end{definizione}

\begin{definizione}[Permutazione con ripetizioni]
	Se tra $n$ oggetti ci sono gruppi di $n_1, n_2, \ldots, n_k$ oggetti identici:
	\[
	\frac{n!}{n_1! \cdot n_2! \cdot \ldots \cdot n_k!}
	\]
\end{definizione}

\begin{esempio}
	Permutazioni della parola ``MISSISSIPPI'' (11 lettere: 4 I, 4 S, 2 P, 1 M):
	\[
	\frac{11!}{4! \cdot 4! \cdot 2! \cdot 1!} = 34650
	\]
\end{esempio}

\begin{definizione}[Disposizioni]
	Numero di modi per scegliere $k$ oggetti da $n$ tenendo conto dell'ordine:
	\[
	D(n,k) = \frac{n!}{(n-k)!} = n \times (n-1) \times \cdots \times (n-k+1)
	\]
\end{definizione}

\begin{definizione}[Combinazioni (coefficiente binomiale)]
	Numero di modi per scegliere $k$ oggetti da $n$ \textbf{senza} considerare l'ordine:
	\[
	C(n,k) = \binom{n}{k} = \frac{n!}{k!(n-k)!}
	\]
\end{definizione}

\textbf{Proprietà dei coefficienti binomiali:}
\begin{align*}
	\binom{n}{0} &= \binom{n}{n} = 1 \\
	\binom{n}{k} &= \binom{n}{n-k} \quad \text{(simmetria)} \\
	\binom{n}{k} &= \binom{n-1}{k-1} + \binom{n-1}{k} \quad \text{(ricorsione - Pascal)} \\
	\sum_{k=0}^{n} \binom{n}{k} &= 2^n
\end{align*}

\subsubsection{Probabilità}

\begin{definizione}[Spazio campionario]
	Insieme $\Omega$ di tutti i possibili risultati di un esperimento.
\end{definizione}

\begin{definizione}[Evento]
	Sottoinsieme $A \subseteq \Omega$ dello spazio campionario.
\end{definizione}

\begin{definizione}[Probabilità]
	Funzione $P: \mathcal{P}(\Omega) \to [0,1]$ tale che:
	\begin{itemize}
		\item $P(\Omega) = 1$
		\item $P(A \cup B) = P(A) + P(B)$ se $A \cap B = \emptyset$ (additività)
	\end{itemize}
	
	\textbf{Caso equiprobabile:}
	\[
	P(A) = \frac{|A|}{|\Omega|} = \frac{\text{casi favorevoli}}{\text{casi possibili}}
	\]
\end{definizione}

\textbf{Regole fondamentali:}
\begin{itemize}
	\item $P(\emptyset) = 0$
	\item $P(\overline{A}) = 1 - P(A)$ (complemento)
	\item $P(A \cup B) = P(A) + P(B) - P(A \cap B)$ (unione)
\end{itemize}

\subsubsection{Probabilità condizionata e indipendenza}

\begin{definizione}[Probabilità condizionata]
	La probabilità di $A$ dato $B$ (con $P(B) > 0$):
	\[
	P(A \mid B) = \frac{P(A \cap B)}{P(B)}
	\]
\end{definizione}

\begin{teorema}[Regola del prodotto]
	\[
	P(A \cap B) = P(A \mid B) \cdot P(B) = P(B \mid A) \cdot P(A)
	\]
\end{teorema}

\begin{definizione}[Eventi indipendenti]
	$A$ e $B$ sono indipendenti se:
	\[
	P(A \cap B) = P(A) \cdot P(B)
	\]
	
	Equivalentemente: $P(A \mid B) = P(A)$
\end{definizione}

\begin{teorema}[Teorema di Bayes]
	\[
	P(A \mid B) = \frac{P(B \mid A) \cdot P(A)}{P(B)}
	\]
\end{teorema}

\subsubsection{Variabili casuali e valore atteso}

\begin{definizione}[Variabile casuale]
	Funzione $X: \Omega \to \mathbb{R}$ che assegna un valore numerico a ogni risultato.
\end{definizione}

\begin{definizione}[Valore atteso]
	Media pesata dei valori:
	\[
	E[X] = \sum_{x} x \cdot P(X = x)
	\]
\end{definizione}

\textbf{Proprietà del valore atteso:}
\begin{itemize}
	\item $E[aX + b] = aE[X] + b$ (linearità)
	\item $E[X + Y] = E[X] + E[Y]$ (sempre vera)
	\item $E[XY] = E[X] \cdot E[Y]$ (solo se $X$ e $Y$ indipendenti)
\end{itemize}

\begin{definizione}[Variabile indicatrice]
	Variabile casuale che vale 1 se un evento si verifica, 0 altrimenti.
	\[
	I_A = \begin{cases}
		1 & \text{se } A \text{ si verifica} \\
		0 & \text{altrimenti}
	\end{cases}
	\]
	
	\textbf{Proprietà fondamentale:} $E[I_A] = P(A)$
\end{definizione}

\begin{esempio}[Numero atteso di confronti in QuickSort]
	Usando variabili indicatrici si dimostra che il numero atteso di confronti è $\Theta(n \log n)$.
\end{esempio}

\subsection{Appendice D (Matrici)}

\subsubsection{Definizioni base}

\begin{definizione}[Matrice]
	Una matrice $A$ di dimensione $m \times n$ è una tabella di $m$ righe e $n$ colonne:
	\[
	A = \begin{bmatrix}
		a_{11} & a_{12} & \cdots & a_{1n} \\
		a_{21} & a_{22} & \cdots & a_{2n} \\
		\vdots & \vdots & \ddots & \vdots \\
		a_{m1} & a_{m2} & \cdots & a_{mn}
	\end{bmatrix}
	\]
	
	Notazione: $A = (a_{ij})$ dove $a_{ij}$ è l'elemento in posizione $(i,j)$
\end{definizione}

\textbf{Tipi speciali:}
\begin{itemize}
	\item \textbf{Matrice quadrata:} $m = n$
	\item \textbf{Matrice identità} $I_n$: diagonale di 1, resto di 0
	\[
	I_3 = \begin{bmatrix}
		1 & 0 & 0 \\
		0 & 1 & 0 \\
		0 & 0 & 1
	\end{bmatrix}
	\]
	\item \textbf{Matrice diagonale:} solo la diagonale può essere non-zero
	\item \textbf{Matrice triangolare superiore:} $a_{ij} = 0$ per $i > j$
	\item \textbf{Matrice simmetrica:} $A = A^T$ (vale $a_{ij} = a_{ji}$)
\end{itemize}

\subsubsection{Operazioni fondamentali}

\textbf{1. Somma:} (solo se $A$ e $B$ hanno stesse dimensioni)
\[
(A + B)_{ij} = a_{ij} + b_{ij}
\]
Complessità: $\Theta(mn)$

\textbf{2. Prodotto per scalare:}
\[
(cA)_{ij} = c \cdot a_{ij}
\]
Complessità: $\Theta(mn)$

\textbf{3. Trasposizione:} scambia righe con colonne
\[
(A^T)_{ij} = a_{ji}
\]
Se $A$ è $m \times n$, allora $A^T$ è $n \times m$

Complessità: $\Theta(mn)$

\textbf{4. Prodotto:} $A$ è $m \times n$, $B$ è $n \times p$
\[
(AB)_{ij} = \sum_{k=1}^{n} a_{ik} b_{kj}
\]
$AB$ è $m \times p$

Complessità (algoritmo naive): $\Theta(mnp)$

Per matrici quadrate $n \times n$: $\Theta(n^3)$

\subsubsection{Proprietà del prodotto}

\begin{itemize}
	\item \textbf{Associativa:} $(AB)C = A(BC)$
	\item \textbf{Distributiva:} $A(B + C) = AB + AC$
	\item \textbf{NON commutativa:} in generale $AB \neq BA$
	\item $(AB)^T = B^T A^T$ (inversione ordine)
	\item $IA = AI = A$ (elemento neutro)
\end{itemize}

\subsubsection{Determinante e matrice inversa}

\begin{definizione}[Determinante]
	Per matrice $2 \times 2$:
	\[
	\det \begin{bmatrix}
		a & b \\
		c & d
	\end{bmatrix} = ad - bc
	\]
	
	Per matrici $n \times n$: sviluppo di Laplace (ricorsivo)
\end{definizione}

\textbf{Proprietà:}
\begin{itemize}
	\item $\det(AB) = \det(A) \cdot \det(B)$
	\item $\det(A^T) = \det(A)$
	\item $\det(I) = 1$
	\item $\det(cA) = c^n \det(A)$ per matrice $n \times n$
\end{itemize}

\begin{definizione}[Matrice inversa]
	$A$ è invertibile se esiste $A^{-1}$ tale che:
	\[
	AA^{-1} = A^{-1}A = I
	\]
	
	\textbf{Condizione necessaria e sufficiente:} $\det(A) \neq 0$
\end{definizione}

\textbf{Proprietà della matrice inversa:}
\begin{itemize}
	\item $(A^{-1})^{-1} = A$
	\item $(AB)^{-1} = B^{-1}A^{-1}$ (inversione ordine)
	\item $(A^T)^{-1} = (A^{-1})^T$
\end{itemize}

\subsubsection{Applicazioni agli algoritmi}

\textbf{1. Grafi:} matrice di adiacenza

$A^k[i][j]$ = numero di cammini di lunghezza $k$ da $i$ a $j$

\textbf{2. Programmazione dinamica:} Floyd-Warshall usa prodotto matriciale generalizzato

\textbf{3. Complessità prodotto matriciale:}
\begin{itemize}
	\item Algoritmo standard: $\Theta(n^3)$
	\item Algoritmo di Strassen: $\Theta(n^{\log_2 7}) \approx \Theta(n^{2.81})$
	\item Algoritmi avanzati: $O(n^{2.37})$ (teorici, non pratici)
\end{itemize}

% ----------------------------------------------------------------
\subsection{Esercizi proposti in aula}
	\subsubsection{Confronto tra BST e Min-Heap}
	
	\begin{definizione}[Proprietà BST]
		In un \textbf{albero binario di ricerca} (BST), per ogni nodo $v$ vale:
		\[
		key(\text{sottoalbero sinistro di } v) \;\leq\; key(v) \;\leq\; key(\text{sottoalbero destro di } v)
		\]
		Questa proprietà garantisce un \emph{ordinamento totale e globale} tra tutti i nodi dell'albero.
	\end{definizione}
	
	\begin{definizione}[Proprietà Min-Heap]
		In un \textbf{min-heap}, per ogni nodo $v$ con figli $u_1, u_2$ vale:
		\[
		key(v) \;\leq\; key(u_1), \qquad key(v) \;\leq\; key(u_2)
		\]
		Questa proprietà garantisce solo un \emph{ordinamento parziale e locale} tra padre e figli diretti. Non esiste alcuna relazione d'ordine imposta tra il figlio sinistro e il figlio destro.
	\end{definizione}
	
	\begin{nota}[Differenza chiave]
		Nel BST la struttura permette di navigare efficientemente verso sinistra o destra per cercare un valore; nel min-heap sappiamo solo che la radice è il minimo globale, ma non abbiamo informazioni sull'ordinamento relativo tra i due sottoalberi.
	\end{nota}
	
	\subsection{Elencare ordinatamente le chiavi: BST vs Min-Heap}
	
	\subsubsection{BST: visita in-order in \texorpdfstring{$O(n)$}{O(n)}}
	
	Nel BST è possibile elencare tutte le $n$ chiavi in ordine crescente in tempo $O(n)$ tramite la \textbf{visita in-order} (sinistra $\to$ radice $\to$ destra), che visita ogni nodo esattamente una volta.
	
	\begin{algorithm}[H]
		\caption{\textsc{InOrder}$(v)$}
		\KwIn{Radice $v$ di un BST}
		\KwOut{Chiavi stampate in ordine crescente}
		\If{$v \neq \textsc{Nil}$}{
			\textsc{InOrder}$(\text{figlio sinistro di } v)$\;
			stampa $key(v)$\;
			\textsc{InOrder}$(\text{figlio destro di } v)$\;
		}
	\end{algorithm}
	
	\subsubsection{Min-Heap: impossibile in \texorpdfstring{$O(n)$}{O(n)}}
	
	\begin{hardnote}
		\textbf{Non è possibile} elencare ordinatamente le chiavi di un min-heap in tempo $O(n)$.
	\end{hardnote}
	
	\paragraph{Dimostrazione per assurdo.}
	Supponiamo per assurdo che esista un algoritmo $A$ che, dato un min-heap di $n$ elementi, ne elenca le chiavi in ordine in tempo $O(n)$.
	
	\begin{enumerate}
		\item La costruzione di un min-heap a partire da $n$ elementi arbitrari richiede tempo $O(n)$ (algoritmo di Floyd).
		\item Per ipotesi, $A$ elenca le chiavi ordinate in $O(n)$.
		\item Concatenando le due operazioni otterremmo un algoritmo di \textbf{ordinamento in tempo}:
		\[
		O(n) + O(n) = O(n)
		\]
		\item Ma questo contraddice il \textbf{limite inferiore} $\Omega(n \log n)$ dimostrato per qualsiasi algoritmo di ordinamento basato su confronti.
	\end{enumerate}
	
	Quindi $A$ non può esistere. $\square$
	
	\begin{softnote}
		In parole semplici: se potessimo estrarre tutte le chiavi in ordine da un heap in $O(n)$, avremmo un algoritmo di sorting lineare, cosa dimostrata impossibile nel modello a confronti. La proprietà del min-heap è \emph{troppo debole} per supportare un'enumerazione ordinata efficiente.
	\end{softnote}
	
	\paragraph{Costo effettivo.}
	Estrarre tutte le $n$ chiavi ordinate da un min-heap richiede $n$ operazioni di \textsc{ExtractMin}, ciascuna di costo $O(\log n)$, per un totale di $O(n \log n)$ — esattamente il costo dell'\textbf{HeapSort}.

% ----------------------------------------------------------------
\subsection{Domande d'esame (Appello 2)}

\begin{hardnote}
	Queste sono domande reali dell'esame.\footnotemark
\end{hardnote}
\footnotetext{Se si studiano solo queste l'esame si passa solo se si è colpiti da un fulmine mentre si fa un salto mortale carpiato suonando il flauto}


\subsubsection*{Definizione di algoritmo}
\textbf{D:} Definizione sufficiente di algoritmo?\\
\textbf{R:} Procedura ben definita che produce output da input in tempo finito.

\subsubsection*{Problema computazionale}
\textbf{D:} Cos'è un problema computazionale?\\
\textbf{R:} Relazione input-output desiderata. Il problema dice \textit{cosa},
l'algoritmo dice \textit{come}.

\subsubsection*{Efficienza}
\textbf{D:} Cosa si intende per efficienza?\\
\textbf{R:} Come crescono le risorse al crescere di $n$
(non il tempo in secondi, dipendente dall'hardware).

\subsubsection*{InsertionSort vs MergeSort}
\textbf{D:} InsSort $= 10n^2$, MergeSort $= 32n\log_2 n$. Quando è più
veloce InsSort?\\
\textbf{R:} Per $n$ piccoli: $10n^2 < 32n\log_2 n$ solo inizialmente.

\subsubsection*{\texorpdfstring{$100n^2$ vs $2^n$}{100n2 vs 2n}}
\textbf{D:} Minimo $n$ tale che $100n^2 < 2^n$?\\
\textbf{R:} $n \ge 15$.

\subsubsection*{Invariante di ciclo}
\textbf{D:} Cos'è un'invariante di ciclo?\\
\textbf{R:} Proprietà vera all'inizio di ogni iterazione. Si dimostra per
inizializzazione, conservazione e conclusione.

\subsubsection*{\texorpdfstring{Vale $\max(f,g) = \Theta(f+g)$?}%
	{Vale max(f,g) = Theta(f+g)?}}
\textbf{R:} Sì: $\max(f,g) \le f+g \le 2\max(f,g)$, con $c_1=1$, $c_2=2$.

\subsubsection*{\texorpdfstring{$2^{n+1}$ e $2^{2n}$ vs $O(2^n)$}%
	{2n+1 e 22n vs O(2n)}}
\textbf{R:} $2^{n+1} = O(2^n)$ (costante 2). $2^{2n} = 4^n \neq O(2^n)$
(richiederebbe $2^n \le c$, impossibile).

\subsubsection*{Costo costante in RAM}
\textbf{R:} Quando la parola di memoria ha dimensione fissa (es.\ 64 bit).

\subsubsection*{Prestazioni su istanze note}
\textbf{R:} Precomputare le soluzioni e restituirle direttamente
quando quell'istanza viene ricevuta.

\subsubsection*{\texorpdfstring{$T(n)=T(n-1)+n^c$}{T(n)=T(n-1)+nc}}
\textbf{R:} $T(n) = \Theta(n^{c+1})$.
Dim.: $\sum_{k=1}^n k^c \approx \int_1^n x^c\,dx = \frac{n^{c+1}}{c+1}$.

\subsubsection*{\texorpdfstring{$T(n)=2T(n/3)+1$}{T(n)=2T(n/3)+1}}
\textbf{R:} Master Theorem, caso 1: $T(n) = \Theta(n^{\log_3 2})$
($\log_3 2 \approx 0.63$).

\subsubsection*{Divide et impera: classe di problemi}
\textbf{R:} Problemi divisibili in 2+ sotto-istanze con combinazione
a costo al più polinomiale.

\subsubsection*{Analisi probabilistica}
\textbf{R:} Uso del calcolo della probabilità nell'analisi degli algoritmi
(non solo caso medio).

\subsubsection*{permute-without-identity}
\textbf{R:} No: con $n=3$ produce solo 3 permutazioni invece delle 5 attese.

\subsubsection*{Hire-Assistant: probabilità}
\textbf{R:} Esattamente 1 assunzione: $1/n$. Esattamente $n$ assunzioni:
$1/n!$.

\subsubsection*{Generatore casuale}
\textbf{R:} Valore intero in un intervallo con probabilità uniforme.

\subsubsection*{Algoritmo randomizzato}
\textbf{R:} Il comportamento dipende dall'input \emph{e} da valori casuali.

\subsubsection*{Ricorrenza con frazioni}
\textbf{D:} $T(n)\le T(10n/11)+T(n/11)+O(n)$?\\
\textbf{R:} $T(n) = O(n\log n)$. Ogni livello costa $O(n)$, profondità
$O(\log n)$.

\subsubsection*{Partition su array}
\textbf{D:} Partition su $[13,19,9,5,12,8,7,4,21,2,6,11]$, pivot $A[11]$?\\
\textbf{R:} $[9,5,8,7,4,2,6,\mathbf{11},21,13,19,12]$, pivot in pos.\ 7.

\subsubsection*{Quicksort con elementi uguali}
\textbf{R:} $O(n^2)$: ogni Partition genera partizione $(n{-}1)$ vs $0$.

\subsubsection*{Partition elementi uguali: fix}
\textbf{R:} Restituisce indice massimo. Fix: flag booleano; se tutti uguali,
restituire $\lfloor(p+r)/2\rfloor$.

\subsubsection*{Quicksort non crescente}
\textbf{R:} Invertire il confronto in Partition: da $\le$ a $\ge$.

\subsubsection*{Rand-Select: caso peggiore}
\textbf{R:} Il pivot è sempre il massimo: dimensione scende di 1 a ogni
passo, $O(n^2)$.

\subsubsection*{\texorpdfstring{IV statistica d'ordine}{IV statistica d'ordine}}
\textbf{D:} IV statistica di $[8,4,8,3,4,2,3,2]$?\\
\textbf{R:} Ordinato: $[2,2,3,3,4,4,8,8]$. IV elemento $= 3$.

\subsubsection*{Eliminare duplicati}
\textbf{D:} Eliminare duplicati in $\Theta(n\log n)$?\\
\textbf{R:} Ordinare, poi scorrere linearmente rimuovendo consecutivi uguali.

\subsubsection*{Build-Max-Heap: indice decrescente}
\textbf{R:} Le foglie sono già heap. MaxHeapify richiede figli sistemati:
si parte dal basso.

\subsubsection*{Visita inorder con pila}
\textbf{R:} Visita in ordine (inorder), $O(n)$.
Pattern: push-sinistra, pop-visita-destra.

\subsubsection*{LCS senza tabella b}
\begin{lstlisting}[language={}, numbers=none, frame=none,
	backgroundcolor=\color{white}, basicstyle=\ttfamily\footnotesize]
	PRINT-LCS(c, X, Y, i, j)
	if c[i,j] == 0 then return
	if X[i] == Y[j]
	PRINT-LCS(c, X, Y, i-1, j-1)
	print X[i]
	else if c[i-1,j] >= c[i,j-1]
	PRINT-LCS(c, X, Y, i-1, j)
	else
	PRINT-LCS(c, X, Y, i, j-1)
\end{lstlisting}
Costo $O(m+n)$: ogni chiamata riduce $i$ o $j$ di 1.

\subsubsection*{Cammino minimo vs TSP}
\textbf{R:} Cammino minimo: sorgente $\to$ destinazione, polinomiale
(Dijkstra). TSP: circuito visitando ogni nodo una volta, NP-difficile.

\subsubsection*{Prima e ultima cifra (Java)}
\begin{lstlisting}[basicstyle=\ttfamily\footnotesize]
	int last  = n % 10;
	int first = n / (int) Math.pow(10, (int)(Math.log10(n)));
\end{lstlisting}

\subsubsection*{\texorpdfstring{\texttt{T extends Comparable}}%
	{T extends Comparable}}
\textbf{R:} La ricerca binaria usa confronti: serve \texttt{compareTo()},
disponibile solo se $T$ estende \texttt{Comparable}.

\subsubsection*{Ereditarietà e attributo omonimo}
\textbf{R:} \texttt{super.getIdA()} per \texttt{A.idA},
\texttt{this.getIdA()} per \texttt{B.idA}.
Le variabili private non sono ereditate.

\subsubsection*{Errore in makeCombination}
\textbf{R:} \texttt{values = numbers} copia solo il riferimento locale.
Fix: copia elemento per elemento.

\subsubsection*{Eccezione in try-with-resources}
\textbf{R:} L'eccezione di \texttt{close()} è soppressa (JLS 14.20.3):
diventa \textit{suppressed exception} della prima.

\subsubsection*{Downcasting Java}
\textbf{D:} \texttt{A value = new B(); B value1 = (A) value; value = null;}?\\
\textbf{R:} \texttt{value} $\to$ null; \texttt{value1} punta all'oggetto
\texttt{B} ancora in memoria.

\subsubsection*{Stack con due code}
\noindent\textbf{Push} $i$-esimo: enqueue in ausiliaria, svuota corrente
nell'ausiliaria, scambia. Costo: $O(i)$. Totale $n$ push: $O(n^2)$.\\
\noindent\textbf{Pop}: dequeue. Costo: $O(1)$.

\subsubsection*{Ricerca in matrice destrorsa}
\textbf{R:} Ricerca binaria su ogni riga. Complessità: $O(n\log n)$.
\begin{lstlisting}[basicstyle=\ttfamily\scriptsize]
	public static <T extends Comparable<T>>
	int[] search(T[][] M, T x) {
		int n = M.length;
		for (int i = 0; i < n; i++) {
			int lo = 0, hi = n - 1;
			while (lo <= hi) {
				int mid = (lo + hi) / 2;
				int cmp = M[i][mid].compareTo(x);
				if      (cmp == 0) return new int[]{i, mid};
				else if (cmp <  0) lo = mid + 1;
				else               hi = mid - 1;
			}
		}
		return null;
	}
\end{lstlisting}


% ================================================================
\subsection{Quiz di Teoria -- Raccolta per Settimane}
% ================================================================

\begin{softnote}
	Le seguenti domande provengono dai quiz settimanali del corso.
	La risposta corretta è evidenziata in \textcolor{correctgreen}{\textbf{verde}}.
\end{softnote}

% -------------------------------------------------------
\subsubsection{Settimana 01}
% -------------------------------------------------------

\paragraph{D1.} Cos'\`e un problema computazionale?
\begin{enumerate}[label=\Alph*.,leftmargin=*]
	\item Una relazione input-output che deve essere determinata da un algoritmo.
	\item \textcolor{correctgreen}{\textbf{Una rappresentazione discorsiva di una relazione input-output desiderata che pu\`o essere determinata da un algoritmo.}} \hfill\textcolor{correctgreen}{$\checkmark$}
	\item Una domanda che richiede l'esecuzione di un algoritmo.
	\item Un problema che deve essere risolto da un algoritmo.
\end{enumerate}

\paragraph{D2.} Minimo $n$ tale che $100n^2 < 2^n$?
\begin{enumerate}[label=\Alph*.,leftmargin=*]
	\item Non esiste un metodo per risolvere la diseguaglianza.
	\item Si deve determinare $n$ per il quale $100n^2 > 2^n$. Risultato: $n \ge 15$.
	\item Si deve determinare $n$ per il quale $\frac{100n^2}{2^n} < 0$. Risultato: $n \ge 15$.
	\item \textcolor{correctgreen}{\textbf{Si deve determinare $n$ per il quale $100n^2 < 2^n$. Risultato: $n \ge 15$.}} \hfill\textcolor{correctgreen}{$\checkmark$}
\end{enumerate}

\paragraph{D3.} InsSort $= 10n^2$, MergeSort $= 32n\log_2 n$. Per quali $n$ InsSort esegue meno passi?
\begin{enumerate}[label=\Alph*.,leftmargin=*]
	\item Si deve risolvere $\frac{10n^2}{32n\log_2 n} < 0$.
	\item Si deve risolvere $32n\log_2 n - 10n^2 < 0$.
	\item \textcolor{correctgreen}{\textbf{Si deve risolvere $10n^2 - 32n\log_2 n < 0$.}} \hfill\textcolor{correctgreen}{$\checkmark$}
	\item $n > 11$.
\end{enumerate}

\paragraph{D4.} Definizione sufficiente di algoritmo.
\begin{enumerate}[label=\Alph*.,leftmargin=*]
	\item Procedura di calcolo ben definita che prende un input e genera un valore come output.
	\item L'astrazione di un programma compilato.
	\item Procedura che prende un insieme di valori come input e genera un valore o un insieme di valori come output.
	\item \textcolor{correctgreen}{\textbf{Procedura di calcolo ben definita che prende un insieme di valori come input e genera un valore, o un insieme di valori, come output in un tempo finito.}} \hfill\textcolor{correctgreen}{$\checkmark$}
\end{enumerate}

\paragraph{D5.} Differenza tra percorso minimo e TSP.
\begin{enumerate}[label=\Alph*.,leftmargin=*]
	\item Il primo chiede percorso minimo, il secondo lo stesso visitando tutte le citt\`a una volta.
	\item Il primo chiede percorso minimo, il secondo un circuito da una citt\`a.
	\item Il primo chiede percorso minimo, il secondo un circuito con partenza $=$ arrivo, visitando tutte le citt\`a.
	\item \textcolor{correctgreen}{\textbf{Il primo chiede percorso minimo tra partenza e arrivo; il secondo lo stesso ma partenza deve essere uguale ad arrivo e si devono visitare tutte le citt\`a una volta sola.}} \hfill\textcolor{correctgreen}{$\checkmark$}
\end{enumerate}

\paragraph{D6.} Cosa si intende per efficienza di un algoritmo?
\begin{enumerate}[label=\Alph*.,leftmargin=*]
	\item \textcolor{correctgreen}{\textbf{La capacit\`a di risolvere istanze di problemi al crescere della dimensione di quest'ultime.}} \hfill\textcolor{correctgreen}{$\checkmark$}
	\item Numero di secondi per risolvere un problema.
	\item Il tempo di esecuzione al crescere della dimensione dell'istanza.
	\item La quantit\`a di un'istanza risolvibile nell'unit\`a di tempo.
\end{enumerate}

% -------------------------------------------------------
\subsubsection{Settimana 02}
% -------------------------------------------------------

\paragraph{D1.} Come migliorare il tempo su istanze note?
\begin{enumerate}[label=\Alph*.,leftmargin=*]
	\item Calcolare le soluzioni del caso peggiore e memorizzarle in un file.
	\item \`E impossibile: non si possono determinare tutte le istanze.
	\item \textcolor{correctgreen}{\textbf{Precomputare le soluzioni per le istanze desiderate; quando arriva quella istanza, l'algoritmo ritorna subito la risposta.}} \hfill\textcolor{correctgreen}{$\checkmark$}
	\item Determinare e memorizzare le soluzioni di tutte le istanze possibili.
\end{enumerate}

\paragraph{D2.} Corrette relazioni tra $2^{n+1}$, $2^{2n}$ e $O(2^n)$.
\begin{enumerate}[label=\Alph*.,leftmargin=*]
	\item $2^{n+1} \neq O(2^n)$, $2^{2n} = O(2^n)$.
	\item $2^{n+1} \neq O(2^n)$, $2^{2n} \neq O(2^n)$.
	\item \textcolor{correctgreen}{\textbf{$2^{n+1} = O(2^n)$, $2^{2n} \neq O(2^n)$.}} \hfill\textcolor{correctgreen}{$\checkmark$}
	\item $2^{n+1} = O(2^n)$, $2^{2n} = O(2^n)$.
\end{enumerate}

\paragraph{D3.} Vale $\max(f(n), g(n)) = \Theta(f(n) + g(n))$?
\begin{enumerate}[label=\Alph*.,leftmargin=*]
	\item \textcolor{correctgreen}{\textbf{S\`i: valgono $f+g = \Theta(\max(f,g))$ e la propriet\`a riflessiva di $\Theta$.}} \hfill\textcolor{correctgreen}{$\checkmark$}
	\item No: $f+g$ potrebbe essere minore di $\max(f,g)$.
	\item No: $\max()$ potrebbe restituire valori diversi per diversi $n$.
	\item S\`i: $\max(f,g)$ \`e sempre minore della somma.
\end{enumerate}

\paragraph{D4.} Cos'\`e un'invariante di ciclo?
\begin{enumerate}[label=\Alph*.,leftmargin=*]
	\item \textcolor{correctgreen}{\textbf{Propriet\`a vera all'inizio di ogni iterazione, usata per dimostrare la correttezza al termine del ciclo.}} \hfill\textcolor{correctgreen}{$\checkmark$}
	\item Propriet\`a vera in qualche iterazione, usata per dimostrare la correttezza.
	\item Propriet\`a che non varia mai nel ciclo ma solo al termine.
	\item Propriet\`a vera all'inizio del ciclo che dimostra la correttezza dell'algoritmo.
\end{enumerate}

\paragraph{D5.} In una RAM, il costo di un'istruzione \`e costante se\ldots
\begin{enumerate}[label=\Alph*.,leftmargin=*]
	\item No, anche con parola di memoria fissa alcune operazioni costano proporzionalmente.
	\item \ldots le operazioni sono solo elementari ($+$, $-$, $*$, $/$).
	\item \ldots ogni operando ha dimensione proporzionale alle parole di memoria.
	\item \textcolor{correctgreen}{\textbf{\ldots la parola di memoria degli operandi ha dimensione fissa.}} \hfill\textcolor{correctgreen}{$\checkmark$}
\end{enumerate}

% -------------------------------------------------------
\subsubsection{Settimana 03--04 (Divide et Impera)}
% -------------------------------------------------------

\paragraph{D1.} Soluzione di $T(n) = T(n-1) + n^c$.
\begin{enumerate}[label=\Alph*.,leftmargin=*]
	\item $T(n) = \Theta(n^c)$.
	\item $T(n) \leq n^c$.
	\item \textcolor{correctgreen}{\textbf{$T(n) = \Theta(n^{c+1})$.}} \hfill\textcolor{correctgreen}{$\checkmark$}
	\item $T(n) = \Theta(n \lg n^c)$.
\end{enumerate}

\paragraph{D2.} Soluzione di $T(n) = 2T(n/3) + 1$.
\begin{enumerate}[label=\Alph*.,leftmargin=*]
	\item $T(n) = \Theta(n^{\log_2 3})$.
	\item \textcolor{correctgreen}{\textbf{$T(n) = \Theta(n^{\log_3 2})$.}} \hfill\textcolor{correctgreen}{$\checkmark$}
	\item $T(n) = \Theta(n^{\lg 2})$.
	\item $T(n) \leq n^{0.61}$.
\end{enumerate}

\paragraph{D3.} Classe di problemi a cui si applica divide et impera efficientemente.
\begin{enumerate}[label=\Alph*.,leftmargin=*]
	\item \textcolor{correctgreen}{\textbf{Problemi in cui un'istanza si divide in 2+ sotto-istanze e la combinazione delle soluzioni richiede costo al pi\`u polinomiale.}} \hfill\textcolor{correctgreen}{$\checkmark$}
	\item Problemi in cui la combinazione ha costo costante.
	\item Problemi in cui le sotto-istanze siano il pi\`u disgiunte possibile.
	\item Problemi divisibili in 1+ sotto-istanze con combinazione polinomiale.
\end{enumerate}

\paragraph{D4.} \textsc{FindMaximumSubarray} su $A = [-10, -2, -4, -11, -6]$.
\begin{enumerate}[label=\Alph*.,leftmargin=*]
	\item \textcolor{correctgreen}{\textbf{Tutti negativi: restituisce l'elemento meno negativo, cio\`e $-2$.}} \hfill\textcolor{correctgreen}{$\checkmark$}
	\item L'algoritmo non si applica con tutti valori negativi.
	\item Restituisce $0$.
	\item Restituisce $-2 + (-4) = -6$.
\end{enumerate}

% -------------------------------------------------------
\subsubsection{Settimana 03--04 (Analisi Probabilistica)}
% -------------------------------------------------------

\paragraph{D1.} \textsc{Hire-Assistant}: probabilit\`a di esattamente 1 assunzione? Di esattamente $n$?
\begin{enumerate}[label=\Alph*.,leftmargin=*]
	\item $\frac{1}{n!}$ e $\frac{1}{n!}$.
	\item $\frac{1}{n}$ e $n\frac{1}{n!}$.
	\item \textcolor{correctgreen}{\textbf{$\frac{1}{n}$ e $\frac{1}{n!}$.}} \hfill\textcolor{correctgreen}{$\checkmark$}
	\item $\frac{1}{n}$ e $\frac{1}{n}$.
\end{enumerate}

\paragraph{D2.} Con analisi probabilistica si intende\ldots
\begin{enumerate}[label=\Alph*.,leftmargin=*]
	\item \ldots uso della probabilit\`a solo quando la distribuzione dell'input \`e nota.
	\item \ldots uso della probabilit\`a nell'analisi del caso medio.
	\item \textcolor{correctgreen}{\textbf{\ldots uso del calcolo della probabilit\`a nell'analisi dei problemi.}} \hfill\textcolor{correctgreen}{$\checkmark$}
	\item \ldots uso della probabilit\`a nell'analisi degli algoritmi.
\end{enumerate}

\paragraph{D3.} Un algoritmo \`e randomizzato quando\ldots
\begin{enumerate}[label=\Alph*.,leftmargin=*]
	\item \ldots il comportamento \`e determinato solo da numeri casuali.
	\item \ldots il comportamento non \`e determinabile.
	\item \ldots i numeri casuali servono solo a velocizzare la computazione.
	\item \textcolor{correctgreen}{\textbf{\ldots il comportamento \`e determinato dall'input \emph{e} da valori prodotti da un generatore casuale.}} \hfill\textcolor{correctgreen}{$\checkmark$}
\end{enumerate}

\paragraph{D4.} Un generatore di numeri casuali\ldots
\begin{enumerate}[label=\Alph*.,leftmargin=*]
	\item \ldots produce un valore di probabilit\`a (tra 0 e 1) con distribuzione uniforme.
	\item \ldots produce un intero diverso dai precedenti in un dato intervallo.
	\item \textcolor{correctgreen}{\textbf{\ldots produce un intero in un dato intervallo con probabilit\`a uniforme sull'intervallo.}} \hfill\textcolor{correctgreen}{$\checkmark$}
	\item \ldots produce un intero con probabilit\`a fissata all'attivazione.
\end{enumerate}

\paragraph{D5.} \textsc{Permute-Without-Identity} genera tutte le permutazioni tranne l'identit\`a con probabilit\`a uniforme?
\begin{enumerate}[label=\Alph*.,leftmargin=*]
	\item S\`i: non genera l'identit\`a e le altre sono equiprobabili.
	\item S\`i: non genera l'identit\`a e le altre significative sono equiprobabili.
	\item No: con $n=3$ genera 4 permutazioni invece di 5.
	\item \textcolor{correctgreen}{\textbf{No: con $n=3$ genera solo 2 permutazioni distinte invece di 5.}} \hfill\textcolor{correctgreen}{$\checkmark$}
\end{enumerate}

% -------------------------------------------------------
\subsubsection{Settimana 05 (QuickSort)}
% -------------------------------------------------------

\paragraph{D1.} Tempo di \textsc{QuickSort} con tutti elementi uguali.
\begin{enumerate}[label=\Alph*.,leftmargin=*]
	\item $O(n\log n) + \Theta(n)$.
	\item \textcolor{correctgreen}{\textbf{$O(n^2)$.}} \hfill\textcolor{correctgreen}{$\checkmark$}
	\item $O(n\lg n)$.
	\item $O(n)$.
\end{enumerate}

\paragraph{D2.} \textsc{Partition} con tutti elementi uguali: quale $q$ restituisce? Come fixare?
\begin{enumerate}[label=\Alph*.,leftmargin=*]
	\item \textcolor{correctgreen}{\textbf{Restituisce l'indice massimo. Fix: flag booleano inizializzato a \texttt{true}; se tutti uguali al pivot, restituire $\lfloor(p+r)/2\rfloor$.}} \hfill\textcolor{correctgreen}{$\checkmark$}
	\item Restituisce l'indice massimo. Fix: ciclo aggiuntivo di verifica.
	\item Restituisce l'indice minimo. Fix: flag booleano.
	\item Restituisce l'indice minimo. Fix: ciclo aggiuntivo.
\end{enumerate}

\paragraph{D3.} Come modificare \textsc{QuickSort} per ordinare in senso non crescente?
\begin{enumerate}[label=\Alph*.,leftmargin=*]
	\item Invertire le due chiamate ricorsive.
	\item \textcolor{correctgreen}{\textbf{Invertire il confronto con il pivot in \textsc{Partition}.}} \hfill\textcolor{correctgreen}{$\checkmark$}
	\item Aggiungere un ciclo che inverte gli elementi.
	\item Ripensare completamente \textsc{Partition}.
\end{enumerate}

\paragraph{D4.} Soluzione di $T(n) \leq T\!\left(\tfrac{10}{11}n\right) + T\!\left(\tfrac{1}{11}n\right) + O(n)$.
\begin{enumerate}[label=\Alph*.,leftmargin=*]
	\item $T(n) = O(n\lg n) + O(n)$.
	\item \textcolor{correctgreen}{\textbf{$T(n) = O(n\lg n)$.}} \hfill\textcolor{correctgreen}{$\checkmark$}
	\item $T(n) = O(n^2)$.
	\item $T(n) = \Omega(n\lg n)$.
\end{enumerate}

\paragraph{D5.} Confronti di \textsc{QuickSort} su $A=[1,2,3,4,5,6,7,8]$?
\begin{enumerate}[label=\Alph*.,leftmargin=*]
	\item \textcolor{correctgreen}{\textbf{$7+6+5+4+3+2+1 = 28$.}} \hfill\textcolor{correctgreen}{$\checkmark$}
	\item $8+7+6+5+4+3+2+1 = 36$.
	\item Dipende dall'implementazione.
	\item $2^8 - 1 = 255$.
\end{enumerate}

% -------------------------------------------------------
\subsubsection{Settimana 05--06 (Selezione)}
% -------------------------------------------------------

\paragraph{D1.} \textsc{Select}$'$ con mediana garantisce $T(n)=O(n)$: come funziona?
\begin{enumerate}[label=\Alph*.,leftmargin=*]
	\item \textcolor{correctgreen}{\textbf{Calcola $x=\textsc{Mediana}(A,p,r)$; usa \textsc{Partition} a 3 parti; ricorre su sottoarray di dimensione $\le n/2$. Ricorrenza: $T(n)\le T(n/2)+O(n)=O(n)$.}} \hfill\textcolor{correctgreen}{$\checkmark$}
	\item Usa \textsc{Partition} a 2 parti; ricorrenza $T(n)\le T(n/2)+O(n)=O(n)$.
\end{enumerate}

\paragraph{D2.} Eliminazione duplicati in $\Theta(n\lg n)$: metodo corretto?
\begin{enumerate}[label=\Alph*.,leftmargin=*]
	\item Ordina e scorre con un ciclo \texttt{for}: $\Theta(n\lg n)+\Theta(n)$.
	\item Ordina e usa doppio ciclo \texttt{for}: $\Theta(n\log n)$.
	\item \textcolor{correctgreen}{\textbf{Cerca la mediana ricorsivamente ed elimina $A_v$: $T(n)=2T(n/2)+O(n)=\Theta(n\log n)$.}} \hfill\textcolor{correctgreen}{$\checkmark$}
	\item Cerca la mediana: $T(n)=T(n/2)+O(n)=\Theta(\log n)$.
\end{enumerate}

\paragraph{D3.} \textsc{Partition}$(A,0,7)$ su $A=[1,2,3,4,5,6,7,8]$: quanti scambi?
\begin{enumerate}[label=\Alph*.,leftmargin=*]
	\item 4. \quad\textbf{B.} \textcolor{correctgreen}{\textbf{8.}} \hfill\textcolor{correctgreen}{$\checkmark$} \quad\textbf{C.} 0. \quad\textbf{D.} 1.
\end{enumerate}

\paragraph{D4.} \textsc{RandSelect} per il minimo di $A=[3,2,9,0,7,5,4,8,6,1]$: caso peggiore?
\begin{enumerate}[label=\Alph*.,leftmargin=*]
	\item L'ultimo elemento \`e sempre il pi\`u piccolo.
	\item L'elemento casuale \`e sempre il pi\`u piccolo.
	\item \textcolor{correctgreen}{\textbf{L'elemento casuale \`e sempre il massimo: la procedura ricorre sempre sul sottoarray di $n-1$ elementi.}} \hfill\textcolor{correctgreen}{$\checkmark$}
	\item L'ultimo elemento \`e sempre il massimo.
\end{enumerate}

\paragraph{D5.} IV statistica d'ordine di $A=[8,4,8,3,4,2,3,2]$?
\begin{enumerate}[label=\Alph*.,leftmargin=*]
	\item 2. \quad\textbf{B.} 8. \quad\textbf{C.} \textcolor{correctgreen}{\textbf{3.}} \hfill\textcolor{correctgreen}{$\checkmark$} \quad\textbf{D.} 4.
\end{enumerate}

% -------------------------------------------------------
\subsubsection{Settimana 06--07 (HeapSort)}
% -------------------------------------------------------

\paragraph{D1.} Perch\'e l'indice in \textsc{BuildMaxHeap} decresce da $\lfloor n/2\rfloor$ a $1$?
\begin{enumerate}[label=\Alph*.,leftmargin=*]
	\item Gli elementi con indice $> \lfloor n/2\rfloor$ sono foglie; \textsc{MaxHeapify} fa scendere la radice.
	\item Si fanno meno cicli.
	\item \textcolor{correctgreen}{\textbf{Gli elementi con indice $> \lfloor n/2\rfloor$ sono foglie (gi\`a max-heap). \textsc{MaxHeapify} richiede sottoalberi gi\`a sistemati: si parte dal basso.}} \hfill\textcolor{correctgreen}{$\checkmark$}
	\item Non c'\`e una ragione specifica.
\end{enumerate}

\paragraph{D2.} Dimostrate che il caso peggiore di \textsc{MaxHeapify} \`e $\Omega(\lg n)$.
\begin{enumerate}[label=\Alph*.,leftmargin=*]
	\item Heap con radice $=0$ e tutti gli altri $=0$: il valore scende a sinistra in $\lfloor\lg n\rfloor$ passi.
	\item Heap con radice $=1$ e tutti gli altri $=0$: il valore scende a sinistra.
	\item \textcolor{correctgreen}{\textbf{Heap con radice $=0$ e tutti gli altri $=1$: l'altezza \`e $\lfloor\lg n\rfloor$; la radice deve scendere fino a una foglia, richiedendo $\lfloor\lg n\rfloor$ chiamate. Caso peggiore $\Omega(\lg n)$.}} \hfill\textcolor{correctgreen}{$\checkmark$}
	\item Come sopra ma il valore scende a destra.
\end{enumerate}

\paragraph{D3.} Cos'\`e uno heap binario come struttura dati?
\begin{enumerate}[label=\Alph*.,leftmargin=*]
	\item Struttura dati basata su un array che rappresenta un albero binario quasi completo.
	\item Un array che rappresenta un albero binario quasi completo.
	\item Struttura composta da un array che rappresenta un albero binario \emph{completo} in modo compatto.
	\item \textcolor{correctgreen}{\textbf{Struttura dati basata su un array che rappresenta un albero binario quasi completo in modo compatto.}} \hfill\textcolor{correctgreen}{$\checkmark$}
\end{enumerate}

\paragraph{D4.} $2i > A.\mathit{heapSize}$: cosa significa?
\begin{enumerate}[label=\Alph*.,leftmargin=*]
	\item $i$ non ha figlio destro.
	\item $i$ non ha figlio sinistro; forse ha figlio destro.
	\item \textcolor{correctgreen}{\textbf{$i$ non ha figlio sinistro, quindi nemmeno destro: \`e una foglia.}} \hfill\textcolor{correctgreen}{$\checkmark$}
	\item $i$ non ha padre.
\end{enumerate}

\paragraph{D5.} Min e max di elementi in uno heap di altezza $h$?
\begin{enumerate}[label=\Alph*.,leftmargin=*]
	\item \textcolor{correctgreen}{\textbf{Minimo $2^h$, massimo $2^{h+1}-1$.}} \hfill\textcolor{correctgreen}{$\checkmark$}
	\item Minimo $2^h$, massimo $2^{h+1}$.
	\item Minimo $2^h-1$, massimo $2^{h+1}-1$.
	\item Minimo $2^h-1$, massimo $2^{h+1}$.
\end{enumerate}

% -------------------------------------------------------
\subsubsection{Settimana 06--07 (Ordinamento Lineare)}
% -------------------------------------------------------

\paragraph{D1.} Risposta in $O(1)$ a ``quanti elementi in $[a,b]$'' con preprocessing $\Theta(n+k)$.
\begin{enumerate}[label=\Alph*.,leftmargin=*]
	\item \textsc{CountingSort}, poi $C[a-1]-C[b]$.
	\item \textcolor{correctgreen}{\textbf{\textsc{CountingSort} in $\Theta(k+n)$; risposta: $C[b]-C[a-1]$, con $C[-1]=0$.}} \hfill\textcolor{correctgreen}{$\checkmark$}
	\item \textsc{CountingSort}, poi $C[a]-C[b]$.
	\item \textsc{CountingSort} in $\Theta(n)$; risposta: $C[b]-C[a-1]$.
\end{enumerate}

\paragraph{D2.} Per quali interi \textsc{BucketSort} \`e conveniente?
\begin{enumerate}[label=\Alph*.,leftmargin=*]
	\item \textcolor{correctgreen}{\textbf{Interi a 8 bit quando $n > 50$.}} \hfill\textcolor{correctgreen}{$\checkmark$}
	\item Interi a 32 bit con $n \approx 10000$.
	\item Interi a 16 bit con $n \approx 50$.
	\item Interi a 8 bit con $n < 50$.
\end{enumerate}

\paragraph{D3.} \textsc{CountingSort} con ciclo da $1$ a $A.\mathit{length}$: corretto? Stabile?
\begin{enumerate}[label=\Alph*.,leftmargin=*]
	\item Corretto: elementi in posizioni $C[k]$ fino a $C[k-1]+1$ nell'ordine dell'input.
	\item Corretto: posizioni da $C[k-1]$ a $C[k-2]+1$.
	\item \textcolor{correctgreen}{\textbf{Corretto: posizioni da $C[k-1]+1$ a $C[k]$, ma in ordine \emph{inverso}: non stabile.}} \hfill\textcolor{correctgreen}{$\checkmark$}
	\item Corretto: posizioni da $C[k-1]$ a $C[k-2]+1$ in ordine inverso.
\end{enumerate}

\paragraph{D4.} Profondit\`a minima di una foglia in un albero di decisione per $n$ elementi.
\begin{enumerate}[label=\Alph*.,leftmargin=*]
	\item \textcolor{correctgreen}{\textbf{$n-1$.}} \hfill\textcolor{correctgreen}{$\checkmark$}
	\item $n$. \quad\textbf{C.} $n\lg n$. \quad\textbf{D.} $n-2$.
\end{enumerate}

\paragraph{D5.} Perch\'e il caso peggiore di \textsc{BucketSort} \`e $\Theta(n^2)$? Come portarlo a $O(n\lg n)$?
\begin{enumerate}[label=\Alph*.,leftmargin=*]
	\item \textcolor{correctgreen}{\textbf{Input non uniforme $\Rightarrow$ tutti in un bucket $\Rightarrow$ \textsc{InsertionSort} a $O(n^2)$. Fix: usare \textsc{MergeSort} invece di \textsc{InsertionSort}.}} \hfill\textcolor{correctgreen}{$\checkmark$}
	\item Input uniforme $\Rightarrow$ tutti in un bucket; fix: \textsc{MergeSort}.
	\item Input non uniforme; fix: randomizzare l'input.
	\item Quasi tutti in un bucket; fix: randomizzare.
\end{enumerate}

% -------------------------------------------------------
\subsubsection{Settimana 08--09 (Hash Table)}
% -------------------------------------------------------

\paragraph{D1.} Inserire $5,28,19,15,20,33,12,17,10$ in tabella hash con $m=9$, $h(k)=k\bmod 9$, concatenamento.
\begin{enumerate}[label=\Alph*.,leftmargin=*]
	\item $0{\to}\emptyset$; $1{\to}10$; $2{\to}20$; $3{\to}12$; $5{\to}5$; $6{\to}33$; $8{\to}17$.
	\item $0{\to}19$; $1{\to}10$; $2{\to}20$; $3{\to}12$; $4{\to}28$; $5{\to}5$; $6{\to}33$; $7{\to}15$; $8{\to}17$.
	\item \textcolor{correctgreen}{\textbf{$0{\to}\emptyset$; $1{\to}10{\to}19{\to}28$; $2{\to}20$; $3{\to}12$; $4{\to}\emptyset$; $5{\to}5$; $6{\to}33{\to}15$; $7{\to}\emptyset$; $8{\to}17$.}} \hfill\textcolor{correctgreen}{$\checkmark$}
	\item $1{\to}10{\to}19$; $2{\to}20{\to}19$; $3{\to}12$; $6{\to}33{\to}15$; $8{\to}17$.
\end{enumerate}

\paragraph{D2.} Cardinalit\`a attesa di $\{\{k,l\}\mid k\neq l,\, h(k)=h(l)\}$ con hashing uniforme.
\begin{enumerate}[label=\Alph*.,leftmargin=*]
	\item \textcolor{correctgreen}{\textbf{Definire $X_{kl}=\mathbf{1}\{h(k)=h(l)\}$; $\Pr\{X_{kl}=1\}=1/m$. Per linearit\`a: $E[Y]=\binom{n}{2}\frac{1}{m}=\frac{n(n-1)}{2m}$.}} \hfill\textcolor{correctgreen}{$\checkmark$}
	\item Nessuna risposta \`e corretta.
	\item $\Theta(n(1+\alpha))$.
	\item $\sum_{i=1}^n \alpha(n-i)$.
\end{enumerate}

\paragraph{D3.} Trovare il massimo in una tabella a indirizzamento diretto $T[0..m-1]$.
\begin{enumerate}[label=\Alph*.,leftmargin=*]
	\item \textcolor{correctgreen}{\textbf{Scorrere $T$ dalla fine verso l'inizio; restituire la chiave del primo slot non NULL. Caso peggiore $O(m)$ (nessuna chiave).}} \hfill\textcolor{correctgreen}{$\checkmark$}
	\item Come sopra ma caso peggiore $O(k)$.
	\item Come sopra ma slot nullo restituisce $T[h(k)]$.
	\item Come sopra ma si scorre dall'inizio.
\end{enumerate}

\paragraph{D4.} Con $m=1000$ e $h(k)=\lfloor m(kA\bmod 1)\rfloor$, $A=(\sqrt{5}-1)/2$: calcolare $h(61),\ldots,h(65)$.
\begin{enumerate}[label=\Alph*.,leftmargin=*]
	\item \textcolor{correctgreen}{\textbf{$h(61)=700$, $h(62)=318$, $h(63)=936$, $h(64)=554$, $h(65)=172$.}} \hfill\textcolor{correctgreen}{$\checkmark$}
	\item Nessuna risposta \`e corretta.
	\item $h(61)=3$, $h(62)=1$, $h(63)=4$, $h(64)=2$, $h(65)=0$.
	\item La funzione non \`e ben definita perch\'e non restituisce un intero.
\end{enumerate}

\paragraph{D5.} Liste ordinate nel concatenamento: come cambiano le prestazioni?
\begin{enumerate}[label=\Alph*.,leftmargin=*]
	\item Ricerca con successo invariata; ricerca senza successo invariata; inserimento pi\`u lento; cancellazione $\Theta(1)$.
	\item \textcolor{correctgreen}{\textbf{Ricerca con successo invariata (no ricerca binaria su lista). Ricerca senza successo migliore (ci si ferma alla prima chiave maggiore). Inserimento $O(n_i)$. Cancellazione $\Theta(1)$ con lista doppiamente concatenata.}} \hfill\textcolor{correctgreen}{$\checkmark$}
	\item Nessuna risposta \`e esatta.
	\item Ricerca con successo $\Theta(1+\alpha/2)$; ricerca senza successo $\Theta(1+\alpha)$.
\end{enumerate}

% -------------------------------------------------------
\subsubsection{Settimana 09--10 (Alberi Binari di Ricerca)}
% -------------------------------------------------------

\paragraph{D1.} Versione ricorsiva corretta di \textsc{Tree-Maximum}.
\begin{enumerate}[label=\Alph*.,leftmargin=*]
	\item \texttt{if (x.right != null) \{ Tree-Maximum(x) \} else \{ return x.key \}}
	\item \texttt{if (x.right != null) \{ Tree-Maximum(x.right) \} else \{ return x.key \}} (manca \texttt{return})
	\item \textcolor{correctgreen}{\textbf{Nessuno degli algoritmi presentati \`e corretto (manca \texttt{return} nella chiamata ricorsiva).}} \hfill\textcolor{correctgreen}{$\checkmark$}
	\item \texttt{if (x.left != null) \{ return Tree-Maximum(x.left) \} else \{ return x.key \}} (usa sinistra: trova il minimo, non il massimo)
\end{enumerate}

\paragraph{D2.} BST-sort con $n$ \textsc{Tree-Insert} + inorder: tempi caso peggiore e caso migliore?
\begin{enumerate}[label=\Alph*.,leftmargin=*]
	\item Caso peggiore $\Theta(n^2)$; caso migliore $O(\lg n)$.
	\item Nessuna risposta \`e corretta.
	\item Caso peggiore $\Theta(n)$; caso migliore $O(n\lg n)$.
	\item \textcolor{correctgreen}{\textbf{Caso peggiore $\Theta(n^2)$ (input ordinato $\Rightarrow$ albero di altezza $n$). Caso migliore $O(n\lg n)$ (albero bilanciato).}} \hfill\textcolor{correctgreen}{$\checkmark$}
\end{enumerate}

\paragraph{D3.} Differenza tra propriet\`a BST e min-heap; si pu\`o elencare in $O(n)$ con min-heap?
\begin{enumerate}[label=\Alph*.,leftmargin=*]
	\item S\`i, si pu\`o elencare in ordine in $O(n)$.
	\item \textcolor{correctgreen}{\textbf{BST: ordine globale (sinistro $\le$ nodo $\le$ destro). Min-heap: ordine locale (padre $\le$ figli). La propriet\`a di min-heap \`e insufficiente per elencare in $O(n)$.}} \hfill\textcolor{correctgreen}{$\checkmark$}
	\item Nessuna risposta \`e corretta.
	\item Non in $O(n)$, ma in $O(n\log n)$.
\end{enumerate}

\paragraph{D4.} Claim del Prof. Bunyan: per ogni $a\in A$, $b\in B$, $c\in C$ vale $a\le b\le c$?
\begin{enumerate}[label=\Alph*.,leftmargin=*]
	\item S\`i: l'esempio di ricerca conferma la tesi.
	\item \textcolor{correctgreen}{\textbf{No. Controesempio: radice 10, figlio sinistro 8 con figli 5 e 9. Cercando 5: $B=\{10,8,5\}$, $C=\{9\}$, ma $10 > 9$.}} \hfill\textcolor{correctgreen}{$\checkmark$}
	\item S\`i: nell'esempio la tesi vale.
	\item S\`i: nell'albero con radice 9 la tesi vale.
\end{enumerate}

\paragraph{D5.} Algoritmo non ricorsivo per inorder con pila.

\begin{lstlisting}[language={}, numbers=none, frame=none,
	backgroundcolor=\color{white}, basicstyle=\ttfamily\footnotesize]
	InorderTreeWalk(T) {
		S = pila vuota
		current = T.root
		finito = false
		while (!finito) {
			if (current != null) {
				S.push(current)
				current = current.left
			} else {
				if (!S.isEmpty()) {
					current = S.pop()
					visit(current)
					current = current.right
				} else { finito = true }
			}
		}
	}
\end{lstlisting}
\textcolor{correctgreen}{Pattern: push verso sinistra; pop, visita, poi vai a destra.} \hfill\textcolor{correctgreen}{$\checkmark$}

% -------------------------------------------------------
\subsubsection{Settimana 10--11 (Programmazione Dinamica)}
% -------------------------------------------------------

\paragraph{D1.} \textsc{Modified-Cut-Rod} con costo fisso $c$ per ogni taglio.
\begin{enumerate}[label=\Alph*.,leftmargin=*]
	\item $r[0]=0$; per $j=1\ldots n$: $q=-\infty$; per $i=1\ldots j-1$: $q=\max(q,p[i]+r[j-i]-c)$; $r[j]=q$.
	\item \textcolor{correctgreen}{\textbf{$r[0]=0$; per $j=1\ldots n$: $q=p[j]$ (nessun taglio, nessun costo); per $i=1\ldots j-1$: $q=\max(q,p[i]+r[j-i]-c)$; $r[j]=q$.}} \hfill\textcolor{correctgreen}{$\checkmark$}
	\item Nessuna risposta \`e corretta.
	\item Come la prima ma ciclo interno fino a $j$ incluso.
\end{enumerate}

\paragraph{D2.} Ricostruire LCS in $O(m+n)$ senza tabella $b$.
\begin{enumerate}[label=\Alph*.,leftmargin=*]
	\item Come la risposta corretta ma stampa \texttt{X[i-1]}.
	\item \textcolor{correctgreen}{\textbf{\texttt{Print-LCS(c,X,Y,i,j)}: se $c[i,j]=0$ ritorna; se $X[i]=Y[j]$ ricorre su $(i-1,j-1)$ e stampa $X[i]$; altrimenti segue il massimo tra $c[i-1,j]$ e $c[i,j-1]$.}} \hfill\textcolor{correctgreen}{$\checkmark$}
	\item Come sopra ma stampa \texttt{Y[j]}.
	\item Come sopra ma stampa \texttt{X[i-1]} nelle diramazioni else.
\end{enumerate}

\paragraph{D3.} Perch\'e la memoization non accelera \textsc{MergeSort}?
\begin{enumerate}[label=\Alph*.,leftmargin=*]
	\item La memoization \`e gi\`a usata implicitamente nell'albero di ricorsione.
	\item Al pi\`u 2 chiamate per coppia di indici; la sovrapposizione \`e minima.
	\item \textcolor{correctgreen}{\textbf{Ogni coppia di indici \`e chiamata al pi\`u una volta: nessuna sovrapposizione tra sottoproblemi, niente da riusare.}} \hfill\textcolor{correctgreen}{$\checkmark$}
	\item Ogni coppia \`e chiamata esattamente una volta.
\end{enumerate}

\paragraph{D4.} DP $O(n)$ per Fibonacci. Vertici e archi del grafo dei sottoproblemi?
\begin{enumerate}[label=\Alph*.,leftmargin=*]
	\item $n+1$ vertici; $v_0$ senza archi uscenti; gli altri 2 archi: totale $2n$.
	\item $n$ vertici; ciascuno con 2 archi: totale $2n$.
	\item $n+1$ vertici; tutti con 2 archi: totale $2(n+1)$.
	\item \textcolor{correctgreen}{\textbf{\texttt{fib[0..n]}, $\mathit{fib}[0]=\mathit{fib}[1]=1$, \texttt{for} $i=2$ to $n$: $\mathit{fib}[i]=\mathit{fib}[i-1]+\mathit{fib}[i-2]$. Grafo: $n+1$ vertici; $v_0,v_1$ senza archi; gli altri con 2 archi: totale $2(n-1)$ archi.}} \hfill\textcolor{correctgreen}{$\checkmark$}
\end{enumerate}

\paragraph{D5.} LCS con $2\min(m,n)$ celle; come ridurre a $\min(m,n)$?
\begin{enumerate}[label=\Alph*.,leftmargin=*]
	\item Mantieni celle successive nella stessa riga e precedenti nella riga precedente.
	\item \textcolor{correctgreen}{\textbf{La riga $i$ dipende solo dalla riga $i-1$: bastano 2 righe di lunghezza $\min(m,n)$. Per ridurre a 1 riga: tieni il valore precedente nella stessa riga e $c[i-1,j-1]$ in una variabile temporanea.}} \hfill\textcolor{correctgreen}{$\checkmark$}
	\item Prendi $\min(m,n)$ come numero di righe invece di colonne.
	\item Nessuna risposta \`e esatta.
\end{enumerate}

% -------------------------------------------------------
\subsubsection{Settimana 13 (Grafi: BFS, DFS, Topologico)}
% -------------------------------------------------------

\paragraph{D1.} Un solo bit per il colore dei nodi in BFS: perch\'e basta?
\begin{enumerate}[label=\Alph*.,leftmargin=*]
	\item \textcolor{correctgreen}{\textbf{Il colore grigio \`e solo illustrativo; l'unico controllo esplicito \`e sul bianco. Bastano bianco e nero $\Rightarrow$ 1 bit.}} \hfill\textcolor{correctgreen}{$\checkmark$}
	\item Il nero \`e solo illustrativo; basta controllare grigio e nero.
	\item Grigio e nero sono illustrativi; basta un solo colore.
\end{enumerate}

\paragraph{D2.} Calcolare $G^2$ con liste di adiacenza e con matrice.
\begin{enumerate}[label=\Alph*.,leftmargin=*]
	\item Liste: $O(|E||V|)$; matrice: Strassen $O(|V|^3)$.
	\item \textcolor{correctgreen}{\textbf{Liste: per ogni $u$, per ogni $v\in\text{Adj}[u]$, per ogni $w\in\text{Adj}[v]$, aggiungere $w$ ad \texttt{Adj}$^2[u]$; costo $O(|V||E|)$. Matrice: quadrato con Strassen, costo $O(|V|^{\lg 7})$.}} \hfill\textcolor{correctgreen}{$\checkmark$}
	\item $O(|V|^3)$ per entrambe.
	\item Aggiunge $u$ ad \texttt{Adj}$^2[w]$; costo $O(|V|^3)$.
\end{enumerate}

\paragraph{D3.} BFS con matrice di adiacenza: complessit\`a?
\begin{enumerate}[label=\Alph*.,leftmargin=*]
	\item Invariata: $O(V+E)$.
	\item Nessuna risposta \`e corretta.
	\item $O(V+V^3)$.
	\item \textcolor{correctgreen}{\textbf{Iterare su tutti gli archi costa $O(V^2)$: complessit\`a $O(V+V^2)$.}} \hfill\textcolor{correctgreen}{$\checkmark$}
\end{enumerate}

\paragraph{D4.} Grado uscente e grado entrante con liste di adiacenza di un grafo orientato.
\begin{enumerate}[label=\Alph*.,leftmargin=*]
	\item Nessuna risposta \`e corretta.
	\item Uscente $O(|V|)$; entrante $O(|V|+|E|)$.
	\item Uscente $O(|E|)$; entrante $O(|E|)$.
	\item \textcolor{correctgreen}{\textbf{Uscente $O(|V|+|E|)$ (scorrendo tutte le liste). Entrante $O(|V|+|E|)$ (contando le occorrenze di ogni nodo come destinazione).}} \hfill\textcolor{correctgreen}{$\checkmark$}
\end{enumerate}

\paragraph{D5.} Calcolare $G^T$ con liste di adiacenza e con matrice.
\begin{enumerate}[label=\Alph*.,leftmargin=*]
	\item Liste: per ogni $u$, per ogni $v\in\text{Adj}[u]$, aggiungere $u$ ad \texttt{Adj}$'[v]$; costo $O(|E|)$. Matrice: copia righe come colonne; costo $O(|V|)$.
	\item \textcolor{correctgreen}{\textbf{Liste: costo $O(|V|+|E|)$. Matrice: per $i=1,\ldots,n$ copia la riga $i$ come colonna $i$ nella nuova matrice; costo $O(|V|^2)$.}} \hfill\textcolor{correctgreen}{$\checkmark$}
\end{enumerate}

\section{Raccolta Completa Esercizi e Quiz di Programmazione}

% --- CLASSE SQUARE ---
\subsection{Modellazione Geometrica: Classe Square}
\begin{definizione}
	Realizzare una classe \texttt{Square} che rappresenti un modello di forme quadrate, gestendo la dimensione del lato in modo che non sia mai negativa \cite{213, 166}.
\end{definizione}

\begin{lstlisting}[caption={Implementazione della classe Square}]
	public class Square {
		private double side;
		
		// Costruttore: imposta la lunghezza se >= 0, altrimenti 0
		public Square(double length) {
			this.side = (length >= 0) ? length : 0;
		}
		
		// Modifica la dimensione del lato solo se il valore e' significativo
		public void modify(double length) {
			if (length >= 0) {
				this.side = length;
			}
		}
		
		public double getArea() {
			return this.side * this.side;
		}
	}
\end{lstlisting}

\newpage
% --- CREA ARRAY ---
\subsection{Manipolazione Array: Duplicazione Consecutiva}
\begin{osservazione}
	Scrivere un metodo statico \texttt{creaArray} che generi $n$ copie consecutive di un array dato. Se $n \le 1$ o l'array è nullo, restituire l'originale senza usare \texttt{java.util.Arrays} \cite{217, 167}.
\end{osservazione}

\begin{lstlisting}[caption={Metodo creaArray}]
	static int[] creaArray(int[] a, int n) {
		if (n <= 1 || a == null) {
			return a;
		}
		int[] arr = new int[a.length * n];
		int index = 0;
		for (int i = 0; i < n; i++) {
			for (int j = 0; j < a.length; j++) {
				arr[index++] = a[j];
			}
		}
		return arr;
	}
\end{lstlisting}

% --- MAX RUN ---
\subsection{Analisi di Sequenze: Lunghezza Massima Serie (Run)}
\begin{nozione}
	Una "serie" è una sequenza di valori adiacenti ripetuti. Il metodo deve restituire la lunghezza della serie più lunga \cite{219, 168}.
\end{nozione}

\begin{lstlisting}[caption={Metodo maxRun}]
	int maxRun(int[] a) {
		if (a == null || a.length == 0) return 0;
		int max = 1; 
		int tmp = 1;
		for (int i = 1; i < a.length; i++) {
			if (a[i] == a[i - 1]) {
				tmp++;
			} else {
				max = Math.max(max, tmp);
				tmp = 1;
			}
		}
		return Math.max(max, tmp);
	}
\end{lstlisting}

% --- EQUALS ---
\subsection{Uguaglianza tra Oggetti: Metodo Equals}
\begin{hardnote}
	Per sovrascrivere correttamente \texttt{Object.equals}, il parametro deve essere di tipo \texttt{Object} e si deve usare \texttt{instanceof} prima del cast \cite{222, 170}.
\end{hardnote}

\begin{lstlisting}[caption={Override di equals per la classe A}]
	public class A {
		private int id;
		private String colore;
		
		public boolean equals(Object o) {
			if (!(o instanceof A a)) {
				return false;
			}
			return this.id == a.id;
		}
	}
\end{lstlisting}

% --- FIBONACCI ---
\subsection{Successione di Fibonacci}
\begin{nota}
	L'approccio iterativo evita l'overhead della ricorsione e gestisce il limite del tipo \texttt{long} (fino a $n=92$) \cite{227, 171}.
\end{nota}

\begin{lstlisting}[caption={Fibonacci iterativo}]
	public static long fib(long i) {
		if (i < 0) return -1;
		if (i == 0) return 0;
		if (i == 1) return 1;
		if (i > 92) return -1;
		
		long prev2 = 0, prev1 = 1, current = 0;
		for (int n = 2; n <= i; n++) {
			current = prev1 + prev2;
			prev2 = prev1;
			prev1 = current;
		}
		return current;
	}
\end{lstlisting}

\newpage
% --- STRINGHE ALTERNATE ---
\subsection{Trasformazione Stringhe: Alterna Upper/Lower}
\begin{lstlisting}[caption={Metodo alternaUpperLower}]
	public String alternaUpperLower(String s) {
		if (s == null || s.isEmpty()) return "";
		StringBuilder sb = new StringBuilder();
		for (int i = 0; i < s.length(); i++) {
			char c = s.charAt(i);
			sb.append(i % 2 == 0 ? Character.toUpperCase(c) : Character.toLowerCase(c));
		}
		return sb.toString();
	}
\end{lstlisting}

% --- INTERFACCE E COMPARABLE ---
\subsection{Interfacce, Immutabilità e Ordinamento}
\begin{esempio}
	Classe \texttt{A} immutabile che implementa un'interfaccia personalizzata \texttt{I} e \texttt{Comparable} \cite{234, 177}.
\end{esempio}

\begin{lstlisting}[caption={Classe A con ordinamento basato su pow()}]
	public class A implements I, Comparable<A> {
		int i, j; // Accesso package richiesto
		
		public A(int i, int j) {
			this.i = i;
			this.j = j;
		}
		
		@Override
		public long pow() {
			return (long) Math.pow(i, j);
		}
		
		@Override
		public int compareTo(A o) {
			return Long.compare(this.pow(), o.pow());
		}
	}
\end{lstlisting}

\newpage
% --- QUICK SORT PARTITION ---
\subsection{Algoritmi di Ordinamento: Partition}
\begin{lstlisting}[caption={Procedura Partition classica}]
	public static int partition(int[] a, int low, int high) {
		if (a == null || low > high) return -1;
		int pivot = a[high];
		int i = low - 1;
		for (int j = low; j < high; j++) {
			if (a[j] <= pivot) {
				i++;
				scambia(a, i, j);
			}
		}
		scambia(a, i + 1, high);
		return i + 1;
	}
\end{lstlisting}

% --- BST SEARCH ---
\subsection{Strutture Dati: Ricerca in un BST}
\begin{lstlisting}[caption={Metodo search in BinarySearchTree}]
	public BinaryNode<T> search(T item) {
		if (item == null) throw new NullPointerException();
		BinaryNode<T> node = getRoot();
		while (node != null) {
			int res = item.compareTo(node.getValue());
			if (res == 0) return node;
			node = (res < 0) ? node.getLeftSon() : node.getRightSon();
		}
		return null;
	}
\end{lstlisting}

% --- NOTE FINALI ---
\begin{nota}
	In caso di collisioni in una tabella hash gestita con concatenamento, la ricerca con successo ha un costo medio di $\Theta(1+\alpha)$ \cite{8, 31}.
\end{nota}

\begin{hardnote}
	Se in un \texttt{try-with-resources} sia il metodo \texttt{close()} che il corpo del \texttt{try} lanciano un'eccezione, quella del \texttt{close()} viene soppressa \cite{35, 58}.
\end{hardnote}